%%
%% Appendix: Key code snippets
%% Independent file, included by ../main.tex
%%

\section{Key Code Snippets}
\label{app:code}

We provide the three most relevant fragments of the released
implementation for reproducibility: the lightweight GT encoder
(Listing~\ref{lst:gt_encoder}), the conditional posterior path inside
the probabilistic model (Listing~\ref{lst:posterior}), and the released
configuration used for the reported results
(Listing~\ref{lst:config}).

\subsection{GT Encoder}
\label{app:code:gt}

The GT encoder is a lightweight 3D CNN that maps the ground-truth volume
to a feature volume and mean-pools it into a 32-dimensional global
descriptor $h_Y$. It is instantiated only in \texttt{gt\_cvae} mode and
only consumed during training.

\begin{lstlisting}[language=Python, caption={models/gt\_encoder.py: the
GT encoder.}, label={lst:gt_encoder}]
class GTEncoder(nn.Module):
    """Maps a GT CT volume to a feature volume, then pools a global
    feature for the variational posterior."""

    def __init__(self, in_ch=1, out_ch=32, mid_ch=16):
        super().__init__()
        self.out_ch = out_ch
        self.net = nn.Sequential(
            nn.Conv3d(in_ch, mid_ch, kernel_size=3, padding=1),
            nn.GroupNorm(4, mid_ch),
            nn.SiLU(inplace=True),
            nn.Conv3d(mid_ch, mid_ch, kernel_size=3, stride=2, padding=1),
            nn.GroupNorm(4, mid_ch),
            nn.SiLU(inplace=True),
            nn.Conv3d(mid_ch, out_ch, kernel_size=3, stride=2, padding=1),
            nn.GroupNorm(4, out_ch),
            nn.SiLU(inplace=True),
        )

    def forward(self, gt_volume, anchor_xyz=None):
        """gt_volume: [B, 1, D, H, W]  ->  feat_volume + gt_global_feat."""
        if gt_volume.dim() == 4:
            gt_volume = gt_volume.unsqueeze(1)
        feat_volume = self.net(gt_volume)
        gt_global_feat = feat_volume.mean(dim=(2, 3, 4))   # [B, 32]
        return {'feat_volume': feat_volume,
                'gt_global_feat': gt_global_feat}
\end{lstlisting}

\subsection{Conditional Posterior Path}
\label{app:code:posterior}

During training the sparse context is encoded once into the pooled
representation $\psi(C)$; in \texttt{gt\_cvae} mode the posterior head
instead receives $\psi(C)\oplus h_Y$ computed from the ground-truth
volume. The decoder always renders from \emph{context} features (the
released \texttt{posterior\_decoder\_views='context'}), so the FiLM
modulated latent is the only privileged channel.

\begin{lstlisting}[language=Python, caption={models/dense\_probability.py:
context encoding (top) and the conditional posterior path (middle), with
FiLM modulation of the decoder features (bottom).},
label={lst:posterior}]
def encode_context(self, data, prepare=True):
    clean = context_inputs(data)
    if prepare:
        clean = self.prepare_views(clean)
    features = self.encode_projs(clean)
    pooled = torch.cat([features['anchor_feats'].mean(-1),
                        observation_descriptor(clean)], -1)
    return ContextState(features, clean, pooled, grid=None)

def infer_posterior(self, data, context, dense=None):
    """GT-conditioned CVAE (gt_cvae): condition q(z | C, Y) on the GT
    volume, not on dense views. GT encoder is never called when
    gt_volume is None."""
    if self.mode == 'dense_np':
        dense = dense if dense is not None else self.encode_dense(data)
        pooled, grid = dense.pooled, dense.grid     # 48 dense views
    else:
        if data.get('gt_volume') is None:
            raise ValueError('gt_cvae posterior requires gt_volume')
        gt = self.gt_encoder(data['gt_volume'],
                             context.data['points_gs'])
        pooled = torch.cat([context.pooled,
                            gt['gt_global_feat']], -1)   # 128+7+32
        grid = None                                     # use_local=False
    return self.posterior_head(pooled), grid

# in sample_state(): draw z and modulate the *context* features
mu, lv = posterior if posterior is not None else prior
z_g = mu if mean else reparameterize(mu, lv, generator)
features['gs_params']['feats'] = self.feature_film(
    context.features['gs_params']['feats'], z_g)
\end{lstlisting}

\subsection{Released Configuration}
\label{app:code:config}

The released configuration (which produced the results reported in
Section~\ref{sec:exp}) fixes the mode to \texttt{gt\_cvae}, uses a
64-dimensional global latent with no local latent
(\texttt{use\_local: false}), a per-dimension free-bits floor of
$0.5$ nats, an annealed $\beta_f=2{\times}10^{-2}$, and the
\emph{prior mean} Monte-Carlo estimator with $N{=}8$ samples at
inference.

\begin{lstlisting}[caption={configs/experiment\_plan/stage1/
prob-gt-cvae.yaml (excerpt): the released model, objective and
evaluation blocks. \texttt{posterior\_decoder\_views} is not a config
key; the code default is \texttt{'context'}.},
label={lst:config}]
model:
  image_encoder:
    out_ch: 128
  anchor_fusion: geometry_transformer
  point_fusion: max
  multiview_transformer:
    d_model: 128
    num_heads: 4
    num_blocks: 2
    feature_scales: [4, 8, 16]
    samples_per_scale: 5
    spatial_window: 3        # window of +-1 anchor neighbors
    token_init: max          # max-fused anchor tokens (ablation: mean)
  gs:
    t: 512
    res: 12
  probabilistic:
    mode: gt_cvae            # GT-conditioned CVAE posterior
    objective: joint
    normalization: groupnorm_anchor_instance
    film_init_std: 1.0e-3    # zero-mean FiLM init (deterministic start)
    train_samples: 1
    latent:
      global_dim: 64         # z in R^64, FiLM-injected
      use_local: false       # global latent only
      sigma_min: 0.001
      sigma_max: 3.0
      initial_sigma: 0.5
    reconstruction:
      likelihood: mse
      observation_sigma: 0.1   # sigma_obs for the Gaussian likelihood
      reduction: mean_query
    kl:
      beta_final: 2.0e-02      # annealed over warmup
      warmup_steps: 500
      free_bits: 0.5           # per-dimension KL floor (nats)
      free_bits_order: per_case
      reduction: sum_latent_mean_batch
    evaluation:
      estimator: prior_mc_mean  # z ~ p(z|C), mean over samples
      num_samples: 8
dataset:
  paired_views:
    dense_count: 48             # training-only side information
    context_counts: [6]         # deployment view budget
\end{lstlisting}